Denoising Diffusion Probabilistic Models (DDPMs) have rapidly become one of the leading families of
generative models, offering a stability of training that Generative Adversarial Networks (GANs) have
historically struggled to achieve. Within a DDPM the diffusion process is fixed and the only learned
component is the network that predicts the noise added at each step, so the choice of that network
governs the quality of the images produced. This thesis asks whether combining residual learning
with the UNet backbone conventionally used for that purpose yields a measurable improvement.

A residual UNet (ResUNet) was implemented in PyTorch, in which every encoder and decoder stage is
constructed from residual blocks rather than plain convolutional stacks, and trained directly on the
standard DDPM objective. It is assessed against two baselines under matched conditions. The first
removes the identity shortcut from each block and is otherwise identical, so that the two networks
have exactly the same parameter count and the comparison isolates residual learning as the single
variable. The second is the denoising architecture of the original DDPM, with four resolution stages
and self-attention, reduced in width so that its parameter count agrees with the proposed network to
within ten per cent. All three were trained on CIFAR-10 at $32 \times 32$ and on CelebA at $64
\times 64$, sharing the diffusion schedule, the optimiser, the number of epochs and the random seed,
and sampling from the same initial noise.

The three networks were compared on restoration accuracy, measured by structural similarity and peak
signal-to-noise ratio against held-out images at six noise levels; on distribution-level quality,
measured by the Fr\'echet Inception Distance and the Inception Score over 2048 generated samples;
and on a nearest-neighbour check against the training set that found no evidence of memorisation on
either dataset.

Every configuration was trained under three random seeds. The residual shortcut lowers the
first-epoch loss in all six paired runs, by five to eleven per cent on average, but that advantage
disappears by the fifth epoch, and at convergence the restoration measures cannot separate the two
networks: the training loss leans towards the plain variant and structural similarity towards the
proposed one, both by less than half of one per cent. Generative quality separates them, on the
deeper configuration only. On CelebA the proposed network leads in every seed, by a mean of $63.31$
FID against a seed-to-seed standard deviation of about $23$; on CIFAR-10 the paired difference is
$-3.46 \pm 13.78$ with the sign changing between seeds. Two networks that denoise equally well are
therefore shown to generate images of very different quality, a discrepancy no restoration measure
anticipated. The reference architecture remains the strongest overall, leading every restoration
measure and CIFAR-10 FID at roughly $2.3$ times the training cost, and is matched only on CelebA FID
by the proposed network.

The hypothesis accordingly receives partial and conditional support: residual learning does not
improve generated image quality at two resolution stages, where it makes no detectable difference,
but at three stages it leads in every seed by a margin larger than the observed spread. Negative Log-Likelihood was excluded on the
grounds that this implementation fixes the reverse-process variances rather than learning them, and
the comparison against Variational Autoencoder and GAN baselines was not carried out.
